% Course project presentation: rubric sections (a)--(g)
% Base paper: Varshney et al., PETS 2025 (VFU-KD / VFU-GA) — align with Part_1.ipynb & Part_2.ipynb
% Optional figures: place PNGs in ./presentation_figures/ (same as evfu_presentation.tex)
\documentclass[11pt,aspectratio=169]{beamer}
\usetheme{Madrid}
\setbeamertemplate{navigation symbols}{}
\setbeamertemplate{footline}[frame number]
\setbeamersize{text margin left=0.9cm,text margin right=0.9cm}

\usepackage{amsmath,amssymb}
\usepackage{booktabs}
\usepackage{graphicx}
\graphicspath{{./presentation_figures/}{./}}

\title{Verification Framework for Vertical Federated Unlearning (VFU)}
\subtitle{Building on PETS 2025 (Varshney et al.) --- E-VFU-style multi-attack auditing}
\author{}
\institute{ML Project --- \texttt{Part\_1.ipynb}, \texttt{Part\_2.ipynb}}
\date{}

\begin{document}

\begin{frame}
  \titlepage
  \vspace{-0.5em}
  \footnotesize\textit{Target length: $\sim$8 min talk + 2 min Q\&A. This deck maps one-to-one to required rubric items (a)--(g).}
\end{frame}

%==============================================================================
\section*{(a) Problem definition}
\begin{frame}{(a) Problem definition}
  \begin{itemize}
    \item \textbf{Regulatory need:} Right to erasure (e.g., GDPR) requires removing influence of withdrawn data from \emph{trained models}, not only from databases.
    \item \textbf{Vertical FL (VFL):} Parties hold \textbf{different features} for the \textbf{same sample IDs}; only embeddings (or gradients) are shared. Unlearning may mean removing a party, a feature set, or specific samples.
    \item \textbf{VFU methods (PETS 2025):} Varshney et al.\ propose \textbf{VFU-KD} (knowledge distillation) and \textbf{VFU-GA} (gradient ascent) for client / feature / sample unlearning in VFL.
    \item \textbf{Our focus:} A \textbf{verification framework}: how do we \emph{audit} whether unlearning actually removed the target signal, beyond naive accuracy checks?
    \item \textbf{Gap:} Output-only membership inference (MIA) is a weak auditor; gradients, forensic probes, and model-level tests can reveal residual leakage even when logits look ``fine.''
  \end{itemize}
\end{frame}

%==============================================================================
\section*{(b) Dataset \& statistics}
\begin{frame}{(b) Dataset and analysis (statistics)}
  \textbf{Sources implemented in the notebooks (after cleaning / encoding):}
  \vspace{0.3em}
  \begin{center}
    \small
    \begin{tabular}{@{}lrrp{5.2cm}@{}}
      \toprule
      \textbf{Dataset} & \textbf{Samples} & \textbf{Features} & \textbf{Target / imbalance} \\
      \midrule
      UCI Adult & 30{,}162 & 14 & Income $>$50K: 7{,}508 positive vs.\ 22{,}654 negative ($\sim$25\% minority) \\
      German Credit & 1{,}000 & 20 & Good credit: 700 vs.\ bad: 300 (70/30) \\
      MIMIC-III & --- & --- & Full pipeline reserved; demo uses 46K synthetic placeholder if no credentialed path \\
      \bottomrule
    \end{tabular}
  \end{center}
  \vspace{0.4em}
  \begin{itemize}
    \item Categorical columns label-encoded; Adult rows with missing values dropped.
    \item \textbf{VFL split:} feature columns partitioned across 2 (or 3) parties; one \textbf{active} party holds labels (binary classification).
    \item Features standardized per party; train/validation/test splits with fixed seed (e.g., 42) for reproducibility.
  \end{itemize}
\end{frame}

%==============================================================================
\section*{(c) Proposed solution --- architecture}
\begin{frame}{(c) Proposed solution: architecture}
  \textbf{Phase 0 --- VFL + unlearning (PETS-style, as in code)}
  \begin{itemize}
    \item \textbf{Bottom models:} Each passive party: MLP encoder maps local features $\to$ embedding.
    \item \textbf{Top model (active party):} Concatenated embeddings $\to$ MLP $\to$ binary logit; trained with \texttt{BCELoss}.
    \item \textbf{VFU-KD (passive party / feature removal):} New student modules trained with \textbf{prediction loss + distillation} from frozen teacher (see loss slide).
    \item \textbf{VFU-GA (sample unlearning):} Short gradient-ascent phases on target samples to suppress their influence (fast but may hurt utility).
  \end{itemize}
  \vspace{0.35em}
  \textbf{Phase 1 --- Verification (E-VFU-style, \texttt{Part\_1.ipynb})}
  \begin{itemize}
    \item \textbf{(1)} Output-based MIA (baseline, as in PETS evaluation).
    \item \textbf{(2)} Gradient-based MIA: membership classifier on gradient norms / layer stats / spectral cues (Random Forest).
    \item \textbf{(3)} FUMA-style forensic checks: canary samples, prediction divergence.
    \item \textbf{(4)} Shadow models + meta-classifier: does the unlearned model resemble ``trained on retain-only'' vs.\ ``full data''?
  \end{itemize}
\end{frame}

%==============================================================================
\section*{(d) Loss function}
\begin{frame}{(d) Loss function and intuition}
  \textbf{Standard VFL training (active top)}
  \[
    \mathcal{L}_{\text{train}} = \mathrm{BCE}(\hat{y}, y).
  \]
  \textbf{VFU-KD passive-party unlearning (\texttt{Part\_2.ipynb})}
  \begin{itemize}
    \item Student top $\theta_a'$ sees embeddings \textbf{without} the removed party; teacher is the original top model.
    \item Per batch:
    \[
      \mathcal{L} = (1-\lambda)\,\underbrace{\mathrm{BCE}(y_{\text{student}}, y)}_{\mathcal{L}_{\text{pred}}}
      + \lambda\,\underbrace{\mathrm{KL}(p_{\text{student}}\ \|\ p_{\text{teacher}})}_{\mathcal{L}_{\text{dist}}}.
    \]
    \item Default $\lambda = 0.3$: balance \textbf{fitting true labels} on the retain set with \textbf{mimicking the old decision boundary} so utility is preserved while the removed party's features are absent.
    \item \textbf{Feature unlearning} on passive side: KL between student and teacher \textbf{embeddings} (distribution match in representation space).
    \item \textbf{VFU-GA intuition:} Ascend on loss for target samples to ``unmemorize'' them; fast but coarser than KD.
  \end{itemize}
\end{frame}

%==============================================================================
\section*{(e) Performance report}
\begin{frame}{(e) Performance report}
  \textbf{Representative numbers from notebook runs (UCI Adult, may vary slightly with rerun)}
  \begin{itemize}
    \item After VFL training ($\sim$50 epochs): accuracy $\approx$ \textbf{0.85}, F1 $\approx$ \textbf{0.69--0.70}, ROC-AUC $\approx$ \textbf{0.91}.
    \item Pre-unlearning evaluation: e.g.\ Acc $\approx$ 0.858, F1 $\approx$ 0.693, AUC $\approx$ 0.912 (logged in \texttt{Part\_2.ipynb}).
    \item \textbf{VFU-KD unlearning:} Total loss $\approx$ 0.25 after 30 epochs with $\mathcal{L}_{\text{pred}}\approx 0.34$, $\mathcal{L}_{\text{dist}}\approx 0.04$ (illustrative run); bar charts compare metrics before vs.\ after unlearning.
    \item \textbf{Verification metrics:} Report AUC of MIA ($\to$ 0.5 desired), FUMA detection rate ($\to$ low), shadow-model probability of ``full-data'' model ($\to$ $<0.5$), and combined pass/fail dashboard (\texttt{presentation\_figures/} if generated).
  \end{itemize}
  \vspace{0.3em}
  \footnotesize Replace with your latest run outputs when presenting; keep one slide as the ``headline'' metrics and move plots to backup slides if needed.
\end{frame}

%==============================================================================
\section*{(f) Comparison with state of the art}
\begin{frame}{(f) Comparison with state-of-the-art methods}
  \small
  \textbf{Horizontal FL unlearning (context)} --- FedEraser, SISA, FedRecover, KNOT, FRU: focus on \textbf{sample} removal with different storage/retrain trade-offs; not the same data geometry as VFL.

  \vspace{0.35em}
  \textbf{Vertical FL unlearning (closer baseline)}
  \begin{center}
    \begin{tabular}{@{}llp{5.5cm}@{}}
      \toprule
      \textbf{Work} & \textbf{Mechanism} & \textbf{Contrast to our project} \\
      \midrule
      VFU-LR (Deng et al.) & Constraint-based & Efficiency vs.\ retraining; different algorithmic family. \\
      Few-shot label unlearning (Gu et al.) & Mixup + GA & Label-side focus; we follow PETS-style VFU-KD/GA. \\
      \textbf{Varshney et al.\ (PETS 2025)} & \textbf{VFU-KD / VFU-GA} & \textbf{Our base:} we implement and \textbf{extend evaluation} with stronger verification. \\
      Certified VFL unlearning (Wang \& Wang) & Certified removal & Theoretical guarantees; we focus on empirical auditing. \\
      \bottomrule
    \end{tabular}
  \end{center}
  \vspace{0.3em}
  \textbf{Our contribution:} Not a new unlearning rule, but a \textbf{multi-attack verification stack} to stress-test whether PETS-style VFU actually hides membership and forensic traces.
\end{frame}

%==============================================================================
\section*{(g) Challenges \& discussion}
\begin{frame}{(g) Challenges / discussion}
  \begin{itemize}
    \item \textbf{Evaluator strength:} Stronger MIAs (including gradient-based and unlearning-specific attacks) can contradict a ``pass'' from output-only MIA alone --- verification must match the threat model.
    \item \textbf{Utility--privacy trade-off:} VFU-GA is fast but can degrade accuracy; VFU-KD needs historical embeddings / teacher access and more epochs.
    \item \textbf{Correlated features across parties:} Removing one party's column may not remove all signal if correlated proxies exist in other parties' views.
    \item \textbf{Scale \& deployment:} Real VFL adds secure aggregation, async parties, and credentialled data (MIMIC); the notebooks use public UCI data for reproducibility.
    \item \textbf{Future work:} Iterative unlearning sequences, communication cost under SMPC, and tighter alignment with certified unlearning theory.
  \end{itemize}
\end{frame}

\begin{frame}{Thank you}
  \centering
  \Large Questions?
\end{frame}

\end{document}
